% -*- coding: utf-8 -*-
% !TEX program = xelatex
\documentclass[plain,noanswer,medmath]{../../template/jnuexam}
\usepackage{../../template/kysx}

\begin{document}


\examtitle{2007}{3}


\loadproblems{1}{2007P1}%载入试卷一
\loadproblems{2}{2007P2}%载入试卷二


\makepart{选择题}{1～10小题，每小题4分，共40分}


%【类似试卷一第一(1)题】
\begin{problem}
当$x \to 0^+$时，与$\sqrt{x}$等价的无穷小量是\pickout{}
\begin{abcd}
\item $1 - \e^{\sqrt{x}}$．
\item $\ln (1 + \sqrt{x})$．
\item $\sqrt{1^{} + \sqrt{x}} - 1$．
\item $1 - \cos \sqrt{x}$．
\end{abcd}
\end{problem}


\useproblem{1}{1}{4}%【同试卷一第一(4)题】


\useproblem{1}{1}{3}%【同试卷一第一(3)题】


\useproblem{2}{1}{8}%【同试卷二第一(8)题】


\begin{problem}
设某商品的需求函数为$Q = 160 - 2p$，其中$Q$，$p$分别表示需要量和价格，
如果该商品需求弹性的绝对值等于$1$，则商品的价格是\pickout{}
\begin{abcd}
\item $10$．
\item $20$．
\item $30$．
\item $40$．
\end{abcd}
\end{problem}


\useproblem{1}{1}{2}%【同试卷一第一(2)题】


\useproblem{1}{1}{7}%【同试卷一第一(7)题】


\useproblem{1}{1}{8}%【同试卷一第一(8)题】


\useproblem{1}{1}{9}%【同试卷一第一(9)题】


\useproblem{1}{1}{10}%【同试卷一第一(10)题】


\makepart{填空题}{11～16小题，每小题4分，共24分}


\begin{problem}
$\lim_{x \to +\infty} \frac{x^3 + x^2 + 1}{2^x + x^3}(\sin x + \cos x) =$ \fillin{}．
\end{problem}


\useproblem{2}{2}{13}%【同试卷二第二(13)题】


\useproblem{2}{2}{15}%【同试卷二第二(15)题】


\begin{problem}
微分方程$\frac{\dy}{\dx} = \frac{y}{x} - \frac{1}{2}\left(\frac{y}{x} \right)^3$
满足$y\big|_{x = 1} = 1$的特解为$y =$ \fillin{}．
\end{problem}


\useproblem{1}{2}{15}%【同试卷一第二(15)题】


\useproblem{1}{2}{16}%【同试卷一第二(16)题】


\makepart{解答题}{17～24小题，共86分}


\begin{problem}[points=10]
设函数$y = y(x)$由方程$y\ln y - x + y = 0$确定，试判断曲线$y = y(x)$在点$(1,1)$附近的凹凸性．
\end{problem}


\useproblem[points=11]{2}{3}{22}%【同试卷二第三(22)题】


%【类似试卷一第三(19)题】
\begin{problem}[points=11]
设函数$f(x)$，$g(x)$在$\;\left[a,b \right]$上连续，在$(a,b)$内二阶可导且存在相等的最大值，又$f(a)$=$g(a)$，$f(b)$=$g(b)$，证明：\par
\begin{enumerate}
  \item 存在$\eta \in(a,b)$，使得$f(\eta) = g(\eta)$；
  \item 存在$\xi \in(a,b)$，使得$f''(\xi) = g''(\xi)$．
\end{enumerate}
\end{problem}


\begin{problem}[points=10]
将函数$f(x) = \frac{1}{x^2 - 3x - 4}$展开成$x - 1$的幂级数，并指出其收敛区间．
\end{problem}


\useproblem[points=11]{1}{3}{21}%【同试卷一第三(21)题】


\useproblem[points=11]{1}{3}{22}%【同试卷一第三(22)题】


\useproblem[points=11]{1}{3}{23}%【同试卷一第三(23)题】


\useproblem[points=11]{1}{3}{24}%【同试卷一第三(24)题】


\end{document}
